% ------------------------------------------------------------------------------
% 5 Discussion
% ------------------------------------------------------------------------------
\section{Discussion}\label{sec:discussion}

\subsection{Principal findings}

Six analytical components, aimed at one 1980s dataset, converge on one
conclusion and enrich it from different angles. The conclusion:
levamisole plus 5-FU after resection of stage B/C colon carcinoma
delivers a large, robust treatment benefit on disease-free
survival---hazard ratio $0.62$ adjusted, $0.59$--$0.61$ Bayesian, $0.60$
under competing risks---corresponding to a five-year absolute gain of
16.8 percentage points, a number needed to treat of six, and 0.63
disease-free years gained in restricted mean survival time. The
enrichment: the effect is stable across seven pre-specified subgroups
with no credible interaction; a modern group-sequential design would
likely have stopped the trial for efficacy at its first interim analysis
(HR $0.52$, $Z = 3.38$ against a boundary of $2.96$), and Bayesian
monitoring would have signalled benefit at a quarter of the evidence;
covariate adjustment would have bought real power at no inferential
cost; and the largest absolute benefits accrue to the highest-risk
patients through the arithmetic of baseline risk rather than
differential response.

The re-enacted monitoring result deserves particular emphasis because it
connects historical data to contemporary trial ethics. Fifty percent of
the 506 DFS events in this dataset contain essentially all the evidence
needed for the superiority conclusion; the remaining half of the
patients' events---years of follow-up and hundreds of
recurrences---added precision but not decisions. Translated to the 1980s
context in which the 625 patients randomized to observation or
levamisole alone in this very trial continued on an ineffective
strategy, an early stop would have redirected patients to the active
regimen years sooner. The counterfactual is stylized (the historical
investigators lacked today's data-management speed and regulatory
framing), but it illustrates with unusual concreteness why alpha-spending
and beta-spending designs became the default.

\subsection{Methodological lessons}

Several results function as transferable teaching points.
\textbf{Non-collapsibility}, usually a paragraph in a textbook, appears
here as a measured quantity: with trial-calibrated prognostic strength,
the unadjusted Cox estimator targets a marginal hazard ratio attenuated
by $0.033$ on the log scale, so unadjusted and adjusted estimates are
\emph{different estimands} even in a perfectly randomized trial
(Simulation Study A). The RMST analysis supplies the constructive
counterpoint: because the restricted-mean difference is collapsible, its
adjusted and unadjusted estimates coincide (0.57 vs.\ 0.63 years within
a five-year horizon, well within sampling error), making RMST an
attractive companion estimand whenever hazard-ratio interpretations
become delicate \citep{royston2011,uno2014,martinussen2013}.
\textbf{Covariate adjustment} raises power from 72.6\% to 80.9\% in
exactly the scenario regulators care about---modest effects, realistic
sample sizes---which is the quantitative argument behind the FDA/ICH
position that adjustment should be pre-specified and reported even when
unadjusted results match \citep{fda2023,kahan2014}. \textbf{Competing
risks} change the estimand, not necessarily the conclusion: the naive
one-minus-Kaplan--Meier overstates cumulative recurrence by about one
percentage point at eight years, yet the Fine--Gray and cause-specific
readings agree on the treatment effect. \textbf{Honest validation}
matters more than model sophistication: the apparent-to-cross-validated
optimism gap reaches 0.12 in the $C$-index for the most aggressive
learner, while a clinical Cox model loses nothing---a reminder that in
small-to-moderate clinical datasets, disciplined regression frequently
matches ensembles. \textbf{Bayesian monitoring} demonstrates both its
appeal (decisions at 25\% of the information) and its governance
question (posterior probabilities do not self-correct for repeated
looks, so their use in confirmatory settings requires either
calibration to frequentist error rates or a genuinely decision-theoretic
framing).

\subsection{Strengths}

The project's strengths are its integration and its reproducibility.
Every module runs from version-controlled code with fixed seeds,
communicates exclusively through files, and regenerates every number and
figure in this paper via a single \texttt{make} invocation; all 34
tabular outputs of the original pipeline were re-derived during
manuscript preparation and verified to be byte-identical to the
committed versions. Quality-control assertions are embedded in the data
pipeline (record counts, within-patient covariate consistency, endpoint
logic), analyses are pre-specified in structure (primary composite
endpoint, adjustment set, subgroup list), and sensitivity analyses were
run in both directions---towards stronger assumptions (parametric
Weibull Bayesian model) and weaker ones (competing risks, splines,
complete case versus imputation, RMST). The simulation studies were
calibrated to the real data-generating mechanism rather than to
stylized parameters, so their conclusions transfer directly to trials of
this type.

\subsection{Limitations}

Five limitations bound the interpretation. First, this is a
\emph{re-analysis of a single historical trial} from an era of different
supportive care and staging; the quantitative findings (especially
absolute risks) should not be read as contemporary practice, and
levamisole itself was later displaced by superior regimens. Second, the
subgroup and CATE analyses are exploratory: with only 324 events in the
two-arm comparison, interaction tests have limited power, the T-learner's
quartile contrasts rest on 60--100 events each, and no external
validation cohort was available to check the predicted benefit ranking;
the honest claim is methodological demonstration, not clinical
deployment. Third, the Weibull baseline is a parametric choice; although
posterior predictive checks pass and LOO prefers it over the
exponential, a piecewise-exponential or semi-parametric Bayesian
baseline would be a natural robustness extension, as would a
time-varying treatment effect in the Bayesian model given the
covariate-level PH violations detected. Fourth, the re-enacted
monitoring and Bayesian sequential analyses re-use the final data to
reconstruct interim states; the reconstruction censors follow-up at a
common time aligned to the event count and is faithful to the ordering
of the event stream, but the exercise remains a counterfactual
simulation, and early-stopping estimates from such re-enactments are
biased toward larger effects (as the interim HR of $0.52$ versus the
final $0.62$ illustrates). Fifth, the prediction models use only ten
baseline covariates from the original case-report forms; modern depth of
data (genomics, imaging) would shift the accuracy ceiling, though the
overfitting lesson of the apparent-versus-validated contrast would only
sharpen.

A further, more technical caveat concerns the simulated operating
characteristics of the sequential design (Study B): the simulated
rejection probability under the alternative (82.9\%) falls short of the
design-theoretic 90\% because the simulation couples an unadjusted
primary analysis (whose marginal hazard ratio is attenuated by the
strong calibration covariates, per Study A) with a simulation sample
size calibrated to the conditional effect, and because futility
adherence removes some would-be winners. Designs that pair
covariate-adjusted primary analyses with group-sequential monitoring
would recover part of this gap---a combination now explicitly
encouraged by regulators \citep{fda2023} and worth studying in its own
right.

\subsection{Conclusions}

This re-analysis set out to demonstrate what a modern biostatistical
toolkit, applied with discipline to a classic randomized trial, can
extract beyond the original publication---and the answer is: the same
conclusion, reached faster (sequential design re-enactment), quantified
more honestly (imputation, competing risks, absolute effects in both
percentage points and disease-free years), understood more deeply
(simulation-calibrated operating characteristics, Bayesian robustness,
collapsibility made concrete), and personalized in direction
(risk-driven absolute benefit). The accompanying repository delivers the
full pipeline---R and Python, 23 figures, 37 tabular outputs, and this
paper---reproducible with one command, in the hope that it serves both
as a reference implementation and as a teaching resource on how
frequentist, Bayesian, and machine-learning perspectives can be made to
interrogate the same data in complementary ways.
